\section{Discussion: Classical vs Quantum Learning Frontiers}\label{sec:frontiers}

\subsection{Why Classical Neural Networks Overwhelmingly Outperform QNNs on Classical Data}
Our empirical findings across the classical and geometric benchmarks align with foundational principles in learning theory, optimization, and physics:
\begin{enumerate}[leftmargin=*]
  \item \textbf{Activation functions and non-unitary dynamics:} Classical deep networks compose unconstrained, continuous non-linear activations such as $\text{ReLU}(z) = \max(0, z)$. ReLU units partition the Euclidean feature space $\mathbb{R}^d$ into polyhedral linear regions whose boundary complexity grows exponentially with network depth without sacrificing gradient propagation. In stark contrast, closed quantum evolution is strictly unitary: $|\psi'\rangle = U|\psi\rangle$. Unitary transformations are linear isometric operators on the complex Hilbert space $\mathbb{C}^{2^n}$ that preserve inner products and $L_2$ state norms. Non-linearity enters a VQC solely through measurement expectations $\langle \psi | \mathcal{O} | \psi \rangle$ and periodic trigonometric functions generated by Pauli rotations ($e^{-i \theta \sigma}$). This restricts QNNs to smooth, truncated Fourier series, rendering piecewise-sharp decision boundaries costly to approximate without high circuit depth.
  \item \textbf{The power of overparameterization without trainability collapse:} In classical deep learning, overparameterization (networks with millions of parameters) smooths the loss landscape, eliminates spurious local minima, and accelerates gradient descent convergence via the Neural Tangent Kernel (NTK) regime. Conversely, overparameterization in quantum variational circuits induces severe expressivity bottlenecks: as the ansatz approaches a 2-design (universal Haar distribution), the gradient variance vanishes exponentially with register size ($\text{Var}[\partial_\theta \mathcal{L}] \in \mathcal{O}(2^{-n})$), a disaster known as \emph{barren plateaus} [5]. Thus, classical models can scale parameters unrestrictedly, whereas QNNs must navigate a razor's edge between limited expressive power and untrainable barren landscapes.
  \item \textbf{Hardware execution and backpropagation efficiency:} On classical CPUs and GPUs, reverse-mode automatic differentiation backpropagates gradients through millions of parameters in a single forward-backward pass ($\mathcal{O}(1)$ overhead relative to inference). On quantum processors, analytic gradient evaluation relies on the parameter-shift rule:
  \begin{equation}
    \frac{\partial \langle \mathcal{O} \rangle}{\partial \theta_j} = \frac{\langle \mathcal{O} \rangle_{\theta_j + \frac{\pi}{2}} - \langle \mathcal{O} \rangle_{\theta_j - \frac{\pi}{2}}}{2},
  \end{equation}
  demanding $2P$ distinct circuit evaluations for $P$ parameters. In the presence of finite-shot shot noise, obtaining $\epsilon$-accurate gradients requires $\mathcal{O}(P / \epsilon^2)$ shots, introducing debilitating wall-clock and sample complexity for training on quantum hardware.
\end{enumerate}

\subsection{Four Problem Regimes Where Quantum Neural Networks Win}
Despite their disadvantages on generic classical tabular vectors, quantum models hold mathematically proven and empirically validated niches:
\begin{enumerate}[leftmargin=*]
  \item \textbf{Quantum-Native Data (States from Physical Systems):}
  \begin{itemize}[leftmargin=*]
    \item \emph{The Problem:} Learning properties, invariants, or classification labels of quantum states directly produced by quantum sensors, quantum repeaters, condensed-matter simulations, or physical quantum computers.
    \item \emph{Why QNNs Win:} Classical networks suffer an exponential measurement bottleneck: they must first perform quantum state tomography (requiring millions of physical measurements) to convert the quantum density matrix into classical numbers before training can begin. A QNN can accept the quantum state $|\psi_{\text{data}}\rangle$ directly into its register via coherent quantum access ($\mathcal{O}(1)$ state transfer), bypassing tomography entirely [7].
    \item \emph{Examples:} Detecting quantum phase transitions in the transverse-field Ising model (PennyLane TFIM) and topological order in bond-alternating Heisenberg chains (Qiskit Phase of Matter \cite{bermejo2024quantum}).
  \end{itemize}
  \item \textbf{Discrete Logarithm and Cryptographic / Group-Theoretic Distributions:}
  \begin{itemize}[leftmargin=*]
    \item \emph{The Problem:} Classifying data points sampled from algebraic groups related to discrete logarithms, hidden subgroups, or modular exponentiation.
    \item \emph{Why QNNs Win:} Liu et al.\ mathematically proved a rigorous, exponential separation in supervised learning [8]: a quantum kernel classifier can classify data generated by discrete log problems in polynomial time ($\mathcal{O}(\text{poly}(n))$), whereas no classical polynomial-time algorithm can achieve classification accuracy better than random guessing unless standard cryptographic hardness assumptions fail.
  \end{itemize}
  \item \textbf{High-Order Parity and Correlated XOR Problems:}
  \begin{itemize}[leftmargin=*]
    \item \emph{The Problem:} Learning $k$-parity classification functions ($y = x_{i_1} \oplus x_{i_2} \oplus \cdots \oplus x_{i_k}$) in high dimensions.
    \item \emph{Why QNNs Win:} Classical MLPs suffer from flat loss plateaus and exponential sample complexity when learning parity functions without intermediate representation cues. Quantum feature maps with multi-qubit Pauli-$Z$ parity measurements naturally evaluate the parity Hamiltonian in $\mathcal{O}(1)$ quantum circuit depth using entangling gates [9].
  \end{itemize}
  \item \textbf{Quantum Kernels with Hard Classical Evaluation:}
  \begin{itemize}[leftmargin=*]
    \item \emph{The Problem:} Evaluating kernel classifiers with Instantaneous Quantum Polynomial-time (IQP) feature maps $U_\Phi(x)$ (e.g., the Havl\'i\v{c}ek quantum feature map benchmark \cite{havlicek2019supervised}).
    \item \emph{Why QNNs Win:} Estimating the inner product $K(x, x') = |\langle 0 | U_\Phi^\dagger(x) U_\Phi(x') | 0 \rangle|^2$ corresponds to computing transition probabilities of commuting quantum circuits that are provably classically $\#\mathbf{P}$-hard to simulate or sample [10]. If a learning problem's natural geometry aligns with this quantum kernel, classical neural networks cannot efficiently approximate the feature space.
  \end{itemize}
\end{enumerate}

\subsection{Empirical Validation on Quantum State and Phase Classification}
To evaluate model performance in regimes where data originates from physical quantum systems, we evaluated the Transverse-Field Ising Model (PennyLane TFIM) and the bond-alternating Heisenberg model (Qiskit Phase of Matter):
\begin{enumerate}
  \item \textbf{Physical Systems and Phase Transitions:} The 4-spin TFIM chain undergoes a quantum phase transition at $h_c = 1.0$ separating ferromagnetic order from paramagnetic disorder. Similarly, the Heisenberg spin chain features a transition between trivial and symmetry-protected topological (SPT) phases.
  \item \textbf{Empirical Findings:} On the PennyLane TFIM benchmark, the depth-2 QNN achieves \pennylaneTfimQNNdTwoEntAccuracy{} test accuracy with 13 trainable parameters, closely matching the wide classical MLP (\pennylaneTfimMLPAccuracy{}). On the Qiskit Phase of Matter task, the depth-2 QNN attains \qiskitPhaseOfMatterQNNdTwoEntAccuracy{} accuracy, substantially outperforming the parameter-matched compact MLP (\qiskitPhaseOfMatterMLPSmallAccuracy{}). This demonstrates that on quantum-native physical data, parameter-matched QNNs exhibit strong representation capability without suffering optimization collapse.
\end{enumerate}

\section{Conclusion}
This case study conducted a rigorous, reproducible comparison of variational quantum classifiers against classical multilayer perceptrons across five domain-purposeful benchmark datasets organized into three distinct inductive-bias categories:
\begin{enumerate}
  \item \textbf{Quantum State and Phase Classification:} PennyLane Transverse-Field Ising Model (TFIM) and Qiskit Phase of Matter.
  \item \textbf{High-Dimensional / Non-Linear Quantum Kernels:} Havl\'i\v{c}ek ZZ-feature map benchmark.
  \item \textbf{Classical Small-Data \& Low-Sample Benchmarks:} Breast Cancer Wisconsin (Diagnostic) and downsampled MNIST Digits (0 vs.\ 1).
\end{enumerate}

On standard Euclidean tabular distributions, classical MLPs consistently outperformed or matched QNNs, demonstrating superior expressivity, optimization smoothness, and wall-clock execution speed. However, on geometric tasks aligned with quantum inductive biases---such as entangled quantum feature map kernels and physical quantum state phase transitions---the variational quantum classifier effectively matched or exceeded capacity-matched classical networks with exceptionally compact parameter counts ($\le 13$ parameters).

Our analysis elucidates why classical neural networks will continue to dominate generic classical tabular applications, while delineating the foundational frontiers---quantum-native data, group-theoretic kernel separations, high-order parity Hamiltonians, and hard quantum feature maps---where quantum learning models offer unique physical and structural advantages. The entire experiment suite, automated Exploratory Data Analysis (EDA) with detailed column-level schemas, comprehensive test coverage, and automated compilation ensure end-to-end scientific provenance and one-command reproducibility.

\begin{thebibliography}{99}
\raggedright
\small\sloppy
\bibitem{schuldkilloran2019}
M. Schuld and N. Killoran, ``Quantum Machine Learning in Feature Hilbert Spaces,'' \emph{Physical Review Letters}, vol. 122, 040504, 2019. doi: \href{https://doi.org/10.1103/PhysRevLett.122.040504}{10.1103/PhysRevLett.122.040504}.

\bibitem{perezsalinas2020}
A. P\'erez-Salinas, A. Cervera-Lierta, E. Gil-Fuster, and J. I. Latorre, ``Data re-uploading for a universal quantum classifier,'' \emph{Quantum}, vol. 4, p. 226, 2020. doi: \href{https://doi.org/10.22331/q-2020-02-06-226}{10.22331/q-2020-02-06-226}.

\bibitem{bowles2024}
J. Bowles, S. Ahmed, and M. Schuld, ``Better than classical? The subtle art of benchmarking quantum machine learning models,'' arXiv:2403.07059, 2024, revised 2026. Available: \url{https://arxiv.org/abs/2403.07059}.

\bibitem{schuld2021fourier}
M. Schuld, R. Sweke, and J. J. Meyer, ``Effect of data encoding on the expressive power of variational quantum-machine-learning models,'' \emph{Physical Review A}, vol. 103, 032430, 2021. doi: \href{https://doi.org/10.1103/PhysRevA.103.032430}{10.1103/PhysRevA.103.032430}.

\bibitem{mcclean2018barren}
J. R. McClean, S. Boixo, V. N. Smelyanskiy, R. Babbush, and H. Neven, ``Barren plateaus in quantum neural network training landscapes,'' \emph{Nature Communications}, vol. 9, 4812, 2018. doi: \href{https://doi.org/10.1038/s41467-018-07090-4}{10.1038/s41467-018-07090-4}.

\bibitem{havlicek2019supervised}
V. Havl\'i\v{c}ek et al., ``Supervised learning with quantum-enhanced feature spaces,'' \emph{Nature}, vol. 567, pp. 209--212, 2019. doi: \href{https://doi.org/10.1038/s41586-019-0980-2}{10.1038/s41586-019-0980-2}.

\bibitem{bermejo2024quantum}
U. Bermejo et al., ``Quantum convolutional neural networks are (effectively) classically simulatable,'' arXiv:2408.12739, 2024. Available: \url{https://arxiv.org/abs/2408.12739}.

\bibitem{huang2022quantum}
H.-Y. Huang et al., ``Quantum advantage in learning from experiments,'' \emph{Science}, vol. 376, pp. 1182--1186, 2022. doi: \href{https://doi.org/10.1126/science.abn7293}{10.1126/science.abn7293}.

\bibitem{liu2021rigorous}
Y. Liu, S. Arunachalam, and K. Temme, ``A rigorous and robust quantum speed-up in supervised machine learning,'' \emph{Nature Physics}, vol. 17, pp. 1013--1017, 2021. doi: \href{https://doi.org/10.1038/s41567-021-01287-z}{10.1038/s41567-021-01287-z}.

\bibitem{bravyi2018quantum}
S. Bravyi, D. Gosset, and R. K\"onig, ``Quantum advantage with shallow circuits,'' \emph{Science}, vol. 362, pp. 308--311, 2018. doi: \href{https://doi.org/10.1126/science.aar3106}{10.1126/science.aar3106}.

\bibitem{bremner2016average}
M. J. Bremner, A. Montanaro, and D. J. Shepherd, ``Average-case complexity versus exponential quantum speedup for commuting circuits,'' \emph{Physical Review Letters}, vol. 117, 080501, 2016. doi: \href{https://doi.org/10.1103/PhysRevLett.117.080501}{10.1103/PhysRevLett.117.080501}.

\bibitem{pennylane}
V. Bergholm et al., ``PennyLane: Automatic differentiation of hybrid quantum-classical computations,'' arXiv:1811.04968, 2018. Available: \url{https://arxiv.org/abs/1811.04968}.
\end{thebibliography}
